\documentclass{beamer}
\usepackage[utf8]{inputenc}
\usepackage[T1]{fontenc}
\usepackage{graphicx}
\usepackage{booktabs}

% Theme
\usetheme{Madrid}
\usecolortheme{default}

% Information
\title{Clustering}
\subtitle{A Short Introduction}

\author{Meyssem Bouzaien}
\institute{Academic Year: 2025 - 2026}
\date{}

\begin{document}

\frame{\titlepage}

% ==========================================
% WHAT IS CLUSTERING?
% ==========================================

\begin{frame}
\frametitle{What is Clustering?}
\begin{center}
\includegraphics[width=\textwidth]{images/01_concept_clustering.png}
\end{center}
\end{frame}

\begin{frame}
\frametitle{Simple Definition}
\begin{block}{Clustering = Grouping}
Dividing data into groups where the points resemble each other
\end{block}

\vspace{1.5em}
\textbf{Goal:}
\begin{itemize}
    \item Points in the same group → \textbf{similar}
    \item Points in different groups → \textbf{different}
\end{itemize}

\vspace{1.5em}
\begin{alertblock}{No labels!}
The algorithm discovers the hidden groups in the data on its own
\end{alertblock}
\end{frame}

% ==========================================
% K-MEANS
% ==========================================

\begin{frame}
\frametitle{K-Means: The Main Algorithm}
\vspace{-0.2cm}
\begin{center}
\includegraphics[width=0.75\textwidth]{images/02_kmeans_principe.png}
\end{center}
\end{frame}

\begin{frame}
\frametitle{How Does K-Means Work?}
\begin{block}{Principle}
K-Means finds K groups by placing centers
\end{block}

\vspace{1em}
\textbf{The 4 steps:}
\begin{enumerate}
    \item \textbf{Choose K} (desired number of groups)
    \item \textbf{Place K centers} at random
    \item \textbf{Assign} each point to the nearest center
    \item \textbf{Repeat} until the centers stop moving
\end{enumerate}

\vspace{1em}
\begin{exampleblock}{Example}
To segment customers: K=3 → Budget, Mid-range, Premium groups
\end{exampleblock}
\end{frame}

% ==========================================
% CHOOSING K
% ==========================================

\begin{frame}
\frametitle{How Do We Choose K?}
\begin{center}
\includegraphics[width=0.9\textwidth]{images/03_methode_coude.png}
\end{center}
\end{frame}

\begin{frame}
\frametitle{The Elbow Method}
\begin{block}{Problem}
K-Means needs to know K in advance. How do we choose it?
\end{block}

\vspace{1em}
\textbf{Solution: Look for the "elbow"}
\begin{itemize}
    \item Test several values of K (1, 2, 3, ..., 10)
    \item Measure the quality for each K
    \item The plot forms an "elbow"
    \item The elbow indicates the optimal K
\end{itemize}

\vspace{1em}
\begin{alertblock}{Rule}
Choose K at the elbow = good trade-off!
\end{alertblock}
\end{frame}

% ==========================================
% THE 3 ALGORITHMS
% ==========================================

\begin{frame}
\frametitle{The 3 Clustering Algorithms}
\begin{center}
\includegraphics[width=\textwidth]{images/04_comparaison_algorithmes.png}
\end{center}
\end{frame}

\begin{frame}
\frametitle{Comparing the 3 Algorithms}
\begin{table}
\centering
\begin{tabular}{lll}
\toprule
\textbf{Algorithm} & \textbf{Advantage} & \textbf{Use} \\
\midrule
K-Means & Fast and simple & Round groups \\
Hierarchical & No need for K & Small datasets \\
DBSCAN & Complex shapes & Noise/Outliers \\
\bottomrule
\end{tabular}
\end{table}

\vspace{1.5em}
\textbf{Which to choose?}
\begin{itemize}
    \item \textbf{K-Means}: Most of the time (fast, effective)
    \item \textbf{Hierarchical}: When you don't know K
    \item \textbf{DBSCAN}: Irregular shapes or noise present
\end{itemize}
\end{frame}

% ==========================================
% APPLICATIONS
% ==========================================

\begin{frame}
\frametitle{Concrete Applications}
\vspace{-0.2cm}
\begin{center}
\includegraphics[width=0.75\textwidth]{images/05_applications.png}
\end{center}
\end{frame}

\begin{frame}
\frametitle{Examples of Use}
\begin{table}
\centering
\small
\begin{tabular}{ll}
\toprule
\textbf{Domain} & \textbf{Application} \\
\midrule
Marketing & Customer segmentation \\
E-commerce & Product recommendations \\
Documents & Organization by topic \\
Security & Anomaly detection \\
Images & Compression (reducing colors) \\
Biology & Gene groups \\
Social networks & Communities \\
\bottomrule
\end{tabular}
\end{table}

\vspace{1em}
\begin{exampleblock}{Typical Case}
A company with 10,000 customers → 4 segments for targeted marketing
\end{exampleblock}
\end{frame}

% ==========================================
% SIMPLE WORKFLOW
% ==========================================

\begin{frame}
\frametitle{How to Do Clustering? (5 Steps)}
\begin{enumerate}
    \item \textbf{Collect} the data
    \vspace{0.5em}
    
    \item \textbf{Normalize} (bring to the same scale)
    \vspace{0.5em}
    
    \item \textbf{Choose K} (elbow method)
    \vspace{0.5em}
    
    \item \textbf{Apply} K-Means
    \vspace{0.5em}
    
    \item \textbf{Interpret} the groups
\end{enumerate}

\vspace{1em}
\begin{alertblock}{The Most Important Thing}
Interpretation! The groups must make sense for your domain.
\end{alertblock}
\end{frame}

% ==========================================
% KEY POINTS
% ==========================================

\begin{frame}
\frametitle{Key Points to Remember}
\begin{block}{1. UNsupervised Learning}
No labels → we discover hidden structures
\end{block}

\vspace{0.5em}
\begin{block}{2. K-Means = The Main Algorithm}
Fast, simple, effective for most cases
\end{block}

\vspace{0.5em}
\begin{block}{3. Elbow Method}
To find the right number of groups (K)
\end{block}

\vspace{0.5em}
\begin{block}{4. Normalize the Data}
ALWAYS bring the features to the same scale
\end{block}

\vspace{0.5em}
\begin{block}{5. Applications Everywhere}
Marketing, documents, anomalies, compression, etc.
\end{block}
\end{frame}

% ==========================================
% CONCLUSION
% ==========================================

\begin{frame}
\frametitle{In Summary}
\textbf{Clustering is:}
\begin{itemize}
    \item Finding groups in unlabeled data
    \item K-Means: the main tool (fast and effective)
    \item Elbow method: for choosing K
    \item Many applications: customers, documents, etc.
\end{itemize}

\vspace{2em}
\begin{alertblock}{The Essential Takeaway}
Clustering reveals hidden structures to help you better understand your data
\end{alertblock}
\end{frame}


\end{document}
